\documentclass[11pt]{article}

\usepackage[T1]{fontenc}
\usepackage[utf8]{inputenc}
\usepackage{lmodern}
\usepackage{amsmath,amssymb,amsthm,mathtools}
\usepackage[a4paper,margin=28mm]{geometry}
\usepackage{microtype}
\usepackage{xurl}
\usepackage[hidelinks]{hyperref}
\usepackage{enumitem}
\usepackage{booktabs}
\usepackage{array}

\newtheorem{theorem}{Theorem}[section]
\newtheorem{lemma}[theorem]{Lemma}
\newtheorem{proposition}[theorem]{Proposition}
\newtheorem{corollary}[theorem]{Corollary}
\newtheorem{definition}[theorem]{Definition}
\newtheorem{remark}[theorem]{Remark}

\newcommand{\N}{\mathbb N}
\newcommand{\WIN}{\ensuremath{\mathrm{WIN}}}
\newcommand{\LOSS}{\ensuremath{\mathrm{LOSS}}}
\newcommand{\DRAW}{\ensuremath{\mathrm{DRAW}}}
\newcommand{\odd}{\operatorname{odd}}
\newcommand{\src}{\operatorname{src}}
\newcommand{\ct}{\operatorname{ct}}
\newcommand{\asl}{\operatorname{asl}}
\newcommand{\Moves}{\operatorname{moves}}

\title{The Two-Player $3n\mathbin{\pm}1$ Game:\
A Binary Normal Form and an Audit-Hardened Reduction}
\author{Grisha Pochuev\\
  \small\href{mailto:n_854@mail.ru}{n\_854@mail.ru}\\
  \small\url{https://github.com/Grisha-Pochuev/3n-plus-minus-1-game}}
\date{Audit-hardened version 5.0, 10 August 2026}

\begin{document}
\maketitle

\begin{abstract}
We study the two-player $3n\pm1$ game on positive odd integers.  Earlier
versions of this manuscript claimed a complete no-DRAW proof.  External and
adversarial audits found genuine defects in the global routing argument.  The
present version is deliberately claim-safe: it retains the binary normal
form, constant-tail identities, factor-removal lemmas, source embedding, and
the repaired factorful exponent-one entry, but it withdraws the unconditional
global theorem until two explicit routing-provenance obligations are proved.

The first unresolved obligation is the seeded $A/B_2$ reset: after a bounded
$A$-streak, a valuation-two $B$ transfer may return to an obligation at a
numerically unrelated source, and the proof must exhibit a certified strict
proof-tree descendant before the control rank is reset.  The second is the
semantic completeness of the factor/high-return/marked-tail router: every
actual DRAW continuation must be shown to fall into one of the declared
constructors with its finite witness provenance intact.  The formerly used
anchor estimate is not assumed; the concrete stress test
$20\to30\to45\to68$ with $B(68)=25>20$ is included as a permanent negative
regression.  Thus this document is a rigorous reduction and repair dossier,
not a claim that the prize problem has already been completely solved.
\end{abstract}

\tableofcontents

\section{The game and the theorem}

The current state is a positive odd integer $n$. For $n>1$ the player to move
chooses one of
\[
  3n-1,\qquad 3n+1,
\]
and divides the chosen value by the largest possible power of $2$, leaving an
odd integer. If the resulting odd integer is $1$, the player who made the move
wins immediately. Players alternate with perfect information.

A position is called \WIN\ for the player to move if at least one move leads to
\LOSS. It is called \LOSS\ if every legal move leads to \WIN. These two classes
are the inductively generated finite-remoteness classes. Any position not in
either class is called \DRAW.

\paragraph{Target no-DRAW statement.}\label{target:nodraw}
Every positive odd starting position has finite remoteness; equivalently, the
$3n\pm1$ game has no \DRAW\ positions and optimal play reaches $1$ after
finitely many moves.  In this version this is the \emph{target statement},
not an unconditional theorem.  The exact remaining hypotheses needed for the
conditional closure are stated in Section~\ref{sec:remaining}.

The theorem is not a statement about arbitrary play. There are legal cycles,
such as
\[
  5\longrightarrow 7\longrightarrow 5,
\]
under suitable choices by the players.

\section{Binary conjugation}

Write an odd state as $n=2m+1$ and define
\[
  F(m)=\left\lceil\frac{3m}{2}\right\rceil .
\]
For a nonnegative integer $x$, let $R(x)$ be the integer obtained by deleting
the maximal alternating suffix of the binary expansion of $x$. Thus
$R(1001_2)=10_2$ and $R(10101_2)=0$.

\begin{proposition}[First normal form]\label{prop:first-normal}
For $m>0$, the two children in $m$-coordinates are exactly
\[
  F(m),\qquad F(R(m)).
\]
\end{proposition}

\begin{proof}
One has
\[
  3(2m+1)-1=2(3m+1),\qquad
  3(2m+1)+1=2(3m+2).
\]
Exactly one of $3m+1$ and $3m+2$ is odd. In the other expression, each
additional factor of two removed corresponds to crossing one additional
alternating bit at the end of the binary expansion of $m$. The process stops
at the first repeated adjacent bit, which is precisely the deletion defining
$R(m)$. Returning to $m$-coordinates gives the displayed pair.
\end{proof}

Since both children lie in the image of $F$, we represent the state $F(q)$ by
$q$. The conjugated game has terminal state $0$ and moves
\[
  A(q)=\left\lceil\frac{3q}{2}\right\rceil,
  \qquad
  B(q)=R(A(q)).
\]

\begin{proposition}[Strict contraction]\label{prop:B-contracts}
For every $q>0$,
\[
  B(q)<q.
\]
\end{proposition}

\begin{proof}
Deleting a nonempty binary suffix gives
\[
  B(q)\le \left\lfloor\frac{A(q)}2\right\rfloor
       \le \left\lfloor\frac{3q+1}{4}\right\rfloor<q
\]
for $q\ge2$, while $B(1)=0$.
\end{proof}

Thus all unbounded difficulty comes from the expanding branch $A$ and from the
unbounded alternating suffix removed by $B$.

\section{Constant-tail coordinates and coefficient sources}

For a positive odd integer $a$, an exponent $r\ge1$, and a phase
$e\in\{0,1\}$, put
\[
  Q_r^e(a)=a2^r-e.
\]
This is a binary word whose final constant block has length $r$ and consists
of the bit $e$.

\begin{lemma}[Long-tail recurrence]\label{lem:long-tail}
For every positive odd $a$, phase $e$, and $r\ge3$,
\[
  A(Q_r^e(a))=Q_{r-1}^e(3a),\qquad
  B(Q_r^e(a))=Q_{r-2}^e(3a).
\]
Moreover
\[
  B(Q_1^e(a))=B(Q_2^e(a)).
\]
\end{lemma}

\begin{proof}
For $r\ge3$, the expanding child still ends in at least two equal bits; hence
its maximal alternating suffix consists only of the final bit. This gives both
long-tail identities. At the boundary, the binary words of $3a-e$ and $6a-e$
differ by appending one bit that extends the same alternating suffix, so the
same prefix remains after deletion.
\end{proof}

Define the embedded three-free coefficient
\[
  J(s)=2A(s)+1.
\]
Every positive odd coefficient $a$ has a unique factorization
\[
  a=3^kJ(s),\qquad k\ge0.
\]
Indeed, remove all factors of $3$ and call the remaining odd integer $c$.
Then $c\not\equiv0\pmod3$ and $(c-1)/2$ is $0$ or $2$ modulo $3$.
The image of
\[
 A(2t)=3t,\qquad A(2t+1)=3t+2
\]
is exactly the set of nonnegative integers not congruent to $1$ modulo $3$,
and $A$ is injective. Hence there is a unique $s$ with
$A(s)=(c-1)/2$, equivalently $c=J(s)$. Unique factorization supplies the
unique exponent $k$. The integer $s$ is called the coefficient source of
$a$ and is denoted $\src(a)$.

\begin{lemma}[Source boundary]\label{lem:source-boundary}
For every $s>0$ there is a phase
\[
  \alpha(s)=1-\left(\left\lfloor\frac{s}{2}\right\rfloor\bmod2\right)
\]
such that
\[
  \odd(3J(s)+1-2\alpha(s))=J(A(s))
\]
with $2$-adic valuation one, while the opposite phase selects $J(B(s))$ with
valuation at least two.
\end{lemma}

\begin{proof}
Put $u=A(s)$. Then $J(s)=2u+1$ and
\[
  3J(s)+1-2e=2(3u+2-e).
\]
The bracket is odd exactly when $e=1-(u\bmod2)$. For $s=2t$ or $2t+1$ one
has $u\bmod2=t\bmod2$, so this phase is exactly $e=\alpha(s)$. If $u$ is
even, then $e=1$ and $3u+1=J(u)$; if $u$ is odd, then $e=0$ and
$3u+2=J(u)$. Thus the valuation-one child is $J(A(s))$. The opposite sign is
the other original $3n\pm1$ child of the odd state $J(s)$; under the
conjugacy of Proposition~\ref{prop:first-normal} that child is $J(B(s))$.
Since it is not the valuation-one branch, its valuation is at least two.
\end{proof}

Multiplying a constant-tail coefficient by $3$ preserves its coefficient
source. Passing through a factor-free signed boundary with coefficient $J(s)$
therefore exposes exactly one ordinary $A/B$ move on the source.

For a positive state $q$, write $\kappa(q)$ for the odd coefficient in its
canonical constant-tail coordinates.

\begin{lemma}[Coefficient pullback]\label{lem:coeff-pullback}
For every positive odd $w$,
\[
  \kappa(3w)=J(R(w)),\qquad
  \kappa(3w-1)=J(R(2w-1)).
\]
\end{lemma}

\begin{proof}
Let $m\ge1$ be the length of the maximal alternating suffix of the odd word
$w$, and put $z=R(w)$. If $z>0$, maximality says that the leading bit of the
suffix equals the last bit of $z$. Because the suffix ends in $1$, its value is
\[
 t_m=\begin{cases}
 (2^m-1)/3,&m\text{ even},\\
 (2^{m+1}-1)/3,&m\text{ odd}.
 \end{cases}
\]
Thus $w=2^mz+t_m$. If $m$ is even, $z$ is even and
\[
  3w+1=2^m(3z+1)=2^mJ(z).
\]
If $m$ is odd, $z$ is odd and
\[
  3w+1=2^m(3z+2)=2^mJ(z).
\]
If $z=0$, the whole word is alternating; then $m$ is odd and
$3w+1=2^{m+1}$, whose odd part is $J(0)=1$. Hence in all cases
$\kappa(3w)=\odd(3w+1)=J(R(w))$.

For the second identity set $v=2w-1$, which is odd. Since
$3v+1=2(3w-1)$, the first identity applied to $v$ gives
\[
 \kappa(3w-1)=\odd(3w-1)=\odd(3v+1)=J(R(v))=J(R(2w-1)).
\]
\end{proof}


\section{Finite outcomes, witnesses, and ranks}

The terminal state $0$ is \LOSS\ with height $h(0)=0$. A finite \WIN\ state has
height
\[
  h(x)=1+\min\{h(y):y\text{ is a \LOSS\ child of }x\},
\]
and a finite \LOSS\ state has height
\[
  h(x)=1+\max\{h(y):y\text{ is a child of }x\}.
\]
A state whose outcome has been proved but which is being carried only to
justify a later incidence is called a \emph{routing witness}. A
\emph{retained proof token} is a finite state for which the proof has also
established the comparison needed by the rank: it either seeds an empty inner
token set at a one-shot seed transition, or replaces a previously retained
token by certified descendants of strictly smaller height. This distinction is
important: a newly discovered finite state is never silently appended to a
ranked multiset.

\begin{lemma}[Multiset token rank]\label{lem:token-rank}
For a finite multiset $M$ of proof heights define
\[
  \Phi(M)=\mathop{\#}_{h\in M}\omega^h,
\]
the natural ordinal sum of the monomials $\omega^h$. Replacing one token of
height $H$ by any finite multiset of certified descendants of heights strictly
below $H$ strictly decreases $\Phi$.
\end{lemma}

\begin{proof}
In Cantor normal form the coefficient of $\omega^H$ decreases by one, while
all inserted terms have lower exponent. No finite collection of lower terms
can compensate for the loss at exponent $H$.
\end{proof}

Every positive state $q$ has unique canonical constant-tail coordinates
$q=Q_r^e(a)$ with $a$ odd and $r\ge1$. If
$a=3^kJ(s)$, define the \emph{state source} $\rho(q)=s$.

\begin{lemma}[Universal state-source bound]\label{lem:source-bound}
For every $q>0$,
\[
  \rho(q)\le \frac{q-1}{6}.
\]
\end{lemma}

\begin{proof}
Write $q=Q_r^e(3^kJ(s))$. Since $r\ge1$ and $e\le1$,
$q\ge2J(s)-1$. Moreover
$J(s)=2\lceil3s/2\rceil+1\ge3s+1$. Hence
$q\ge6s+1$.
\end{proof}

\begin{lemma}[Boundary coefficient descent]\label{lem:boundary-coeff}
Let $a$ be positive and odd and
\[
 p=B(Q_1^e(a))=B(Q_2^e(a)).
\]
If $p>0$, then its canonical coefficient satisfies
\[
  \kappa(p)\le\frac{3a+1}{4}.
\]
In particular $\kappa(p)<a$ for $a\ge3$.
\end{lemma}

\begin{proof}
The common child is $p=R(3a-e)$. Deleting a nonempty suffix gives
$p\le\lfloor(3a-e)/2\rfloor$. For every positive integer $y$, the odd
coefficient in its canonical constant-tail coordinates is at most $(y+1)/2$.
Therefore
\[
 \kappa(p)\le\frac{p+1}{2}\le\frac{3a+1}{4}.
\]
The final inequality is strict relative to $a$ for $a\ge3$.
\end{proof}

\section{Side identities and typed boundary objects}

The finite routing system uses two elementary binary identities. We record
them here because they are the only places where an exceptional residue class
enters.

\begin{lemma}[Ordinary side relation]\label{lem:side}
Put $S(q)=B(q)$ and $T(q)=B(A(q))$. If
\[
  q\not\equiv1,3,12,14\pmod{16},
\]
then
\[
  T(q)\in\{A(S(q)),B(S(q))\}.
\]
If $q$ is in one of the four exceptional classes, then $A(q)$ is not
exceptional.
\end{lemma}

\begin{proof}
Let $x=A(q)$ and let $k$ be the length of the maximal alternating suffix of
$x$. Direct reduction of $A(q)$ modulo $8$ shows that $k\ge3$ is possible
exactly for $q\equiv1,3,12,14\pmod{16}$. Outside these classes $k=1$ or
$2$. Writing $x$ as $S(q)$ followed by that one- or two-bit alternating word,
the formulas
\[
  A(2u)=3u,\qquad A(2u+1)=3u+2
\]
show that $A(x)$ is obtained from $A(S(q))$ by appending the corresponding
one- or two-bit alternating word. Deleting its maximal alternating suffix
therefore leaves either $A(S(q))$ or $B(S(q))$. The four exceptional residue
classes map under $A$ to $2,5,10,13\pmod{16}$, none of which is exceptional.
\end{proof}

\begin{corollary}[No four consecutive wins]\label{cor:no-four-win}
For no $q$ can all four states
\[
  q,A(q),A^2(q),A^3(q)
\]
be \WIN.
\end{corollary}

\begin{proof}
If $q$ and $A(q)$ are both \WIN, then $B(q)$ must be \LOSS; the same holds at
the next two positions. If $q$ is ordinary, Lemma~\ref{lem:side} makes
$B(A(q))$ a child of the \LOSS\ state $B(q)$, impossible. Hence $q$ must be
exceptional. Applying the same argument to $A(q)$ forces $A(q)$ exceptional,
contrary to Lemma~\ref{lem:side}.
\end{proof}

\begin{lemma}[Alternating lift identity]\label{lem:alt-lift}
For every positive integer $z$ there are $m\ge1$ and
$\delta=1-(z\bmod2)$ such that
\[
  3z+1=Q_m^\delta(J(R(z))).
\]
If $R(z)>0$, one may take $m$ equal to the length of the maximal alternating
suffix of $z$; if $R(z)=0$, one takes one more than the binary length of $z$.
\end{lemma}

\begin{proof}
First suppose $r=R(z)>0$, let $k$ be the alternating-suffix length and put
$\epsilon=r\bmod2$. Then
\[
 z=2^k r+t_{k,\epsilon},\qquad
 t_{k,0}=\Big\lfloor\frac{2^k}{3}\Big\rfloor,
 \quad t_{k,1}=\Big\lfloor\frac{2^{k+1}}{3}\Big\rfloor.
\]
Since $J(r)=3r+1+\epsilon$, direct substitution gives
\[
 3t_{k,\epsilon}+1+\delta=2^k(1+\epsilon),
\]
the two parity cases being exactly the two remainders of a power of $2$
modulo $3$. Hence
\[
 3z+1+\delta=2^kJ(r),
\]
which is the required identity with $m=k$. If $R(z)=0$, the whole binary word
of $z$ is alternating and starts with $1$, so
$z=t_{k,1}$. The same calculation gives
$3z+1+\delta=2^{k+1}=2^{k+1}J(0)$.
\end{proof}

\begin{lemma}[Exceptional side attachment]\label{lem:exceptional-side}
Let $q$ be exceptional and suppose
\[
 q\text{ is \WIN},\qquad A(q)\text{ is \DRAW},\qquad
 \ell=B(q)\text{ is \LOSS}.
\]
Then for some $m\ge1$ and $\delta\in\{0,1\}$ the two children of $A(q)$ are
exactly
\[
  Q_m^\delta(J(\ell)),\qquad Q_{m+1}^\delta(J(\ell)),
\]
and this adjacent pair contains a continuing \DRAW.
\end{lemma}

\begin{proof}
Write $q=16t+c$ with $c\in\{1,3,12,14\}$. The four direct exceptional
formulas are
\[
\begin{array}{c|c|c|c}
c&\ell&B(A(q))&A^2(q)\\ \hline
1&R(6t)&18t+1&36t+3\\
3&R(6t+1)&18t+4&36t+8\\
12&R(6t+4)&18t+13&36t+27\\
14&R(6t+5)&18t+16&36t+32.
\end{array}
\]
In every row put respectively $z=6t,6t+1,6t+4,6t+5$. Then
$\ell=R(z)$ and $B(A(q))=3z+1$. By Lemma~\ref{lem:alt-lift},
\[
  B(A(q))=Q_m^\delta(J(\ell))
\]
for some $m\ge1$. The last column is $2B(A(q))+\delta$, so it equals
$Q_{m+1}^\delta(J(\ell))$. These are precisely the two children of the
\DRAW\ state $A(q)$; therefore neither is \LOSS\ and at least one is \DRAW.
\end{proof}

For a source $x$ and phase $e$, let $O(x,e)$ denote the following exact
\emph{DRAW obligation}:
\begin{enumerate}[label=(\alph*)]
\item $Q_1^e(J(x))$ is \DRAW; or
\item there is a \DRAW\ state whose two children are exactly
      $Q_2^e(J(x))$ and $Q_1^e(J(x))$.
\end{enumerate}
This definition introduces no outcome assumption beyond the displayed
\DRAW\ state.

\begin{lemma}[Hidden parent for a typed three/two frame]\label{lem:hidden-parent}
Let $a$ be positive, odd and not divisible by $3$, and suppose
\[
  a\equiv1+e\pmod3.
\]
Put $H=Q_3^e(a)$, $G=Q_2^e(a)$ and
\[
  K=\frac{2H+e-1}{3}.
\]
Then $K$ is an integer and
\[
  \Moves(K)=\{H,G\}.
\]
\end{lemma}

\begin{proof}
For $e=0$ the congruence gives $H\equiv2\pmod3$ and
$K=(2H-1)/3$; for $e=1$ it gives $H\equiv0\pmod3$ and $K=2H/3$.
These are exactly the two inverse branches of the injective map
$A(q)=\lceil3q/2\rceil$. The last three bits of $H$ are equal, so the
contracting child removes only the last bit and equals $G$.
\end{proof}

\section{Fixed-tail factor removal and entry typing}

The next two lemmas isolate the only raw entry phenomenon that is not already
an obligation or a factor frame.

\begin{lemma}[Fixed-tail factor-removal diamond]\label{lem:fixed-tail-factor}
Let $c$ be positive and odd, $e\in\{0,1\}$, and $r\ge2$. Put
\[
  X=Q_r^e(3c),\qquad G=Q_r^e(c),\qquad H=Q_{r+1}^e(c),\qquad Y=A(G).
\]
Then
\[
  \Moves(H)=\{X,Y\}.
\]
If $X$ is \DRAW, either $G$ is \DRAW, or the diamond determines a finite
routing witness together with its exact incidence to $G,H,X,Y$ and the other
child of $G$.
\end{lemma}

\begin{proof}
Since $r+1\ge3$, Lemma~\ref{lem:long-tail} gives $A(H)=X$. For $r\ge3$ it
also gives $B(H)=Q_{r-1}^e(3c)=A(G)$; for $r=2$ the boundary identity gives
the same equality. Thus $B(H)=Y$.

Because $X$ is \DRAW, $H$ is not \LOSS. If $H$ is \WIN, then $Y$ is a
\LOSS\ child of $H$. Assume $H$ is \DRAW. Then $Y$ is not \LOSS. If $G$ is
\LOSS, then $Y$ is \WIN. If $G$ is \WIN, either $Y$ itself is a \LOSS\
witness, or the other child of $G$ is a \LOSS\ witness. These are all outcome
orientations. In every finite branch the witness and the parent-child
incidence just displayed are known exactly; no token-rank comparison is
asserted at this stage.
\end{proof}

\begin{definition}[Entry controls]\label{def:entry-controls}
A \emph{boundary entry} is an obligation $O(x,e)$ or, more generally, a
factor-free adjacent frame
\[
 \{Q_{m+1}^e(J(x)),Q_m^e(J(x))\},\qquad m\ge1,
\]
known to contain a continuing \DRAW. Its remaining tail length is recorded as
$\tau$. The typed three/two frame satisfying Lemma~\ref{lem:hidden-parent} is
a distinguished $m=2$ boundary constructor.
A \emph{loss-sibling entry} is one of the finitely many local constructors
in which a continuing \DRAW\ is accompanied by a finite \WIN/\LOSS\ routing
witness with a proved parent/child incidence: the fixed-tail diamond of Lemma~\ref{lem:fixed-tail-factor}, the exponent-one predecessor diamond of Lemma~\ref{lem:exp-one-lift}, an ordinary side diamond of Lemma~\ref{lem:side}, the exceptional side attachment of Lemma~\ref{lem:exceptional-side}, or the hidden-parent frame of Lemma~\ref{lem:hidden-parent}. The constructor records its remaining
constant-tail/factor exponent as the local counter $\tau$.
\end{definition}

\begin{lemma}[Entry typing]\label{lem:entry-typing}
Every finite branch produced by Lemma~\ref{lem:fixed-tail-factor} is a
loss-sibling entry in the sense of Definition~\ref{def:entry-controls}; no
ranked multiset is changed in order to make this classification.
\end{lemma}

\begin{proof}
The proof of Lemma~\ref{lem:fixed-tail-factor} gives the finite witness and
its exact incidence to $G,H,X,Y$ and the other child of $G$. These are
precisely the data of the first loss-sibling constructor. Classification is
therefore immediate; subsequent tail/factor normalization belongs to the
typed routing relation and is ranked there by $\tau$ and, when available,
by descendant-token comparisons.
\end{proof}

\begin{lemma}[Exponent-one lift]\label{lem:exp-one-lift}
Let $a$ be positive and odd and $k\ge1$. If
$R_k=Q_1^e(3^ka)$ is \DRAW, then
\[
  P_{k-1}=Q_2^e(3^{k-1}a)
\]
is either \DRAW\ or \WIN; in the latter case its other child is a finite
\LOSS\ witness.
\end{lemma}

\begin{proof}
The boundary recurrence gives $A(P_{k-1})=R_k$. A state with a \DRAW\ child
cannot be \LOSS. If $P_{k-1}$ is \WIN, the \DRAW\ child cannot witness the
win, so its other child is \LOSS.
\end{proof}

\begin{proposition}[Arbitrary factorful exponent-one entry]\label{prop:factorful-entry}
Let $a=J(s)$ and $k>0$. If $Q_1^e(3^ka)$ is \DRAW, then after finitely many
outcome-compatible local steps one obtains either a typed entry of
Lemma~\ref{lem:entry-typing}, or the factor-free exponent-two state
\[
  Q_2^e(a)\text{ \DRAW}.
\]
\end{proposition}

\begin{proof}
Apply Lemma~\ref{lem:exp-one-lift}. Its finite branch is a loss-sibling entry by Definition~\ref{def:entry-controls}. Otherwise $Q_2^e(3^{k-1}a)$ is \DRAW. Apply
Lemma~\ref{lem:fixed-tail-factor} with $r=2$ repeatedly. Whenever its finite
branch occurs, Lemma~\ref{lem:entry-typing} supplies a typed entry. If it never
occurs, one factor of $3$ is removed at the same tail exponent at each step,
so after $k-1$ steps the displayed factor-free exponent-two state is reached.
\end{proof}

\begin{proposition}[Complete exponent-one/two boundary reduction]\label{prop:boundary-reduction}
Let $s\ge0$, $k\ge0$, $e\in\{0,1\}$, and let
\[
  Q_r^e(3^kJ(s)),\qquad r\in\{1,2\},
\]
be \DRAW. For $s>0$, after finitely many local steps at least one of the following occurs: a typed entry is obtained, the retained source strictly
decreases, or the factor-free exponent-two state $Q_2^e(J(s))$ is \DRAW.
For $s=0$ the corresponding statement is Lemma~\ref{lem:zero-source}.
\end{proposition}

\begin{proof}
Assume $s>0$. There are four cases. If $r=1,k=0$, the state is the first
alternative of the obligation $O(s,e)$ and is already a boundary entry. If
$r=1,k>0$, use Proposition~\ref{prop:factorful-entry}. If $r=2,k=0$, we are
at the displayed factor-free exponent-two state. Finally, if $r=2,k>0$, apply
Lemma~\ref{lem:fixed-tail-factor} repeatedly with $r=2$: a finite branch is a
loss-sibling entry by Definition~\ref{def:entry-controls}, while otherwise
one factor of $3$ is removed at the same exponent at each step and after $k$
steps the factor-free exponent-two state is reached. These cases are disjoint
and exhaustive.
\end{proof}

\section{Factor-free exponent-two base entry}

Let $x>0$, $a=J(x)$ and
\[
  P=Q_1^e(a),\quad U=Q_2^e(a),\quad
  b=B(P)=B(U),\quad F=A(U)=Q_1^e(3a),\quad g=1-e.
\]
Assume $U$ is \DRAW\ and $x$ is the least coefficient source of any \DRAW.  This is a \emph{factor-free} coefficient $J(x)$.  The common child $b$ cannot be $0$, because $0$ is \LOSS\ and a \DRAW\ has no \LOSS\ child.  By Lemma~\ref{lem:boundary-coeff}, the canonical coefficient of $b$ is strictly below $J(x)$.  After deleting any factor of $3$ its three-free coefficient is smaller still, hence it equals $J(t)$ with $t<x$ because $J$ is strictly increasing.  Therefore the state source of $b$ is below $x$, and $b$ is not \DRAW.  No such inference is made for a coefficient $3^kJ(x)$ with $k>0$.
As a child of a \DRAW\ it is not \LOSS, so
\[
  b\text{ is \WIN},\qquad F\text{ is \DRAW}.
\]
\begin{lemma}[First-factor anchor identity]\label{lem:first-factor}
For the states above,
\begin{equation}\label{eq:first-factor}
  9J(x)+1-2e=2^vJ(b),\qquad v\ge1.
\end{equation}
\end{lemma}

\begin{proof}
Put $a=J(x)$. If $e=0$, then $P=2a$, $A(P)=3a$ and $b=R(3a)$.
Apply the first identity of Lemma~\ref{lem:coeff-pullback} to the odd integer
$w=3a$:
\[
 \odd(9a+1)=\kappa(9a)=J(R(3a))=J(b).
\]
If $e=1$, then $P=2a-1$, $A(P)=3a-1$ and $b=R(3a-1)$. The second identity
of Lemma~\ref{lem:coeff-pullback}, again with $w=3a$, gives
\[
 \odd(9a-1)=J(R(6a-1)).
\]
By the boundary identity of Lemma~\ref{lem:long-tail},
$R(6a-1)=R(3a-1)=b$. Hence in both phases the odd part of
$9a+1-2e$ is exactly $J(b)$, which proves \eqref{eq:first-factor}.
\end{proof}

\begin{proposition}[Base-entry attachment]\label{prop:base-entry}
Every factor-free exponent-two \DRAW\ at a globally minimal source has a
finite continuation which either strictly lowers the retained source or
produces a boundary/loss-sibling typed entry.
\end{proposition}

\begin{proof}
Suppose first that $e=1-\alpha(x)$. The first source-boundary numerator is
divisible by $4$, while
\[
9J(x)+1-2e=3(3J(x)+1-2e)-2(1-2e)\equiv2\pmod4.
\]
Thus $v=1$. Writing $u=A(x)$ gives $b=3u+1$. The two children of the \DRAW\
state $F$ are
\[
  T=Q_1^g(J(b)),\qquad C=R(T),
\]
and direct appended-bit deletion gives
$C\in\{A(b),B(b)\}$. If $T$ is \DRAW, this is the boundary entry
$O(b,g)$. If $C$ is \DRAW, the other child $\ell$ of the \WIN\ state $b$ is a
finite \LOSS\ witness. If $b$ is ordinary, this is the ordinary-side
loss-sibling constructor. If $b$ is exceptional and $C=A(b)$,
Lemma~\ref{lem:exceptional-side} gives an adjacent \DRAW\ frame over
$\ell$, hence a boundary/loss-sibling entry. It remains only the exceptional
orientation $C=B(b)$. In that orientation
$b\le(9x+5)/2$ and the alternating suffix of $A(b)$ has length at least three,
so
\[
 C=B(b)\le\frac{A(b)}8\le\frac{3b+1}{16}.
\]
Lemma~\ref{lem:source-bound} then yields
\[
 \rho(C)\le\frac{C-1}{6}<\frac{3b+1}{96}
 \le\frac{27x+17}{192}<x,
\]
a strict retained-source exit.

Now suppose $e=\alpha(x)$. Lemma~\ref{lem:source-boundary} gives
$3J(x)+1-2e=2J(A(x))$, hence $v\ge2$. If $v\ge4$, then
\[
16J(b)\le9J(x)+1\le27x+19,
\]
and $J(b)\ge3b+1$ implies
\[
  b\le\frac{9x+1}{16}<x.
\]
Thus the continuing \DRAW\ has smaller retained source. If $v=2$, the two
children of $F$ are exactly $Q_2^g(J(b))$ and $Q_1^g(J(b))$, which is the
second form of $O(b,g)$. If $v=3$, they are the typed three/two frame
$Q_3^g(J(b)),Q_2^g(J(b))$. Reducing
\eqref{eq:first-factor} modulo $3$ gives
\[
  J(b)\equiv1+g\pmod3,
\]
so Lemma~\ref{lem:hidden-parent} applies. The hidden parent is a boundary
entry; its finite side, if one is exposed, is a loss-sibling entry. Hence no
untyped base row remains.
\end{proof}

\section{The source-zero branch}

For $k\ge0$, $r\ge1$ and $e\in\{0,1\}$ put
\[
  S_{k,r}^e=Q_r^e(3^k).
\]
These are precisely the constant-tail states with coefficient source zero.

\begin{lemma}[Zero-source boundary normalization]\label{lem:zero-source}
If a source-zero \DRAW\ exists, then after finitely many outcome-compatible
moves one reaches a typed boundary or loss-sibling entry.  No ranked token is
installed and no control reset is used before this finite normalization is
complete.
\end{lemma}

\begin{proof}
For $r\ge3$, Lemma~\ref{lem:long-tail} gives
\[
  \Moves(S_{k,r}^e)=\{S_{k+1,r-1}^e,S_{k+1,r-2}^e\}.
\]
Following a \DRAW\ child strictly lowers $r$, so a \DRAW\ with $r\in\{1,2\}$
is reached after finitely many moves. If $k=0$ already at that boundary, the
four possibilities are explicit:
\[
 Q_1^1(1)=1\ \WIN,\quad Q_2^1(1)=3\ \WIN,\quad
 Q_1^0(1)=2\ \LOSS,\quad Q_2^0(1)=4\ \LOSS.
\]
Hence a boundary \DRAW\ must have accumulated at least one factor of $3$.

If the reached boundary has tail exponent one, put $n=k+1$.  If it has
tail exponent two, its children are the common raw child and
$S_{k+1,1}^e$.  If the raw child is \DRAW, again put $n=k+1$; otherwise the
raw child is \WIN\ and forces $S_{k+1,1}^e$ to be \DRAW, in which case put
$n=k+2$.  Thus in all cases one obtains a boundary \DRAW\ configuration whose
raw/signed exits are governed, for some $n\ge1$, by
\[
  3^n+1-2e=2^jJ(y).
\]
The exact $2$-adic valuation is
\[
\begin{array}{c|cc}
 &n\text{ odd}&n\text{ even}\\ \hline
 e=0&j=2&j=1\\
 e=1&j=1&j=2+v_2(n).
\end{array}
\]
For $j\ge2$ the raw and signed exits form the adjacent factor-free frame
\[
  Q_{j-1}^{1-e}(J(y)),\qquad Q_j^{1-e}(J(y)),
\]
containing a continuing \DRAW; by Definition~\ref{def:entry-controls} this is a boundary entry, with its finite tail length recorded in $\tau$. For $j=1$ one has
\[
  y=\frac{3^{n-1}-1}{2}.
\]
When $e=0$ the raw exit is $A(y)$; when $e=1$ it is $B(y)$. The raw and
signed states are the two exits of the same boundary \DRAW\ configuration.
If the signed exit is \DRAW, it is already the exponent-one form of an
obligation. If the raw exit is the continuing \DRAW\ and the signed exit is
finite, that finite state is necessarily \WIN; choosing a canonical \LOSS\
child of it supplies an exact finite routing witness, hence a loss-sibling
entry. If both exits are \DRAW, use the signed obligation. Thus every row
reaches a typed entry after a finite pre-entry computation; no impossible
``source below zero'' comparison is used.
\end{proof}

\begin{lemma}[$A$-selecting side diamond]\label{lem:a-side}
Let $x>0$, $e=\alpha(x)$, $w=Q_1^e(J(x))$, and $y=A(x)$. Then
\[
  B(w)\in\{A(y),B(y)\}.
\]
\end{lemma}

\begin{proof}
The source-boundary identity gives
$A(w)=Q_1^{1-e}(J(y))=4A(y)+1+e$. For $e=0$ the appended bits are $01$;
for $e=1$ they are $10$. If the first appended bit repeats the last bit of
$A(y)$, deleting the alternating suffix leaves $A(y)$; otherwise the deletion
continues through the maximal alternating suffix of $A(y)$ and leaves
$B(y)$.
\end{proof}

\begin{lemma}[Universal $B$-transfer diamond]\label{lem:b-diamond}
Let $x>0$, let $e=1-\alpha(x)$, put $a=J(x)$ and $g=1-e$, and factor
\[
  3a+1-2e=2^jJ(y),\qquad j\ge2,
\]
so $y=B(x)$. Define
\[
 P=Q_1^e(a),\quad U=Q_2^e(a),\quad
 D=Q_j^g(J(y)),\quad C=Q_{j-1}^g(J(y)),\quad F=A(U).
\]
Then $D=A(P)$ and $C=B(P)=B(U)$. If $X$ is the common child of $D,C$ and
$Y$ is the other child of $C$, then
\[
  B(F)=Y.
\]
Consequently an obligation in the $B$-selecting phase transfers a continuing
\DRAW\ to the adjacent frame $\{D,C\}$; when $j=2$ the transferred frame is
again an obligation in the opposite phase.
\end{lemma}

\begin{proof}
The formulas for $D$ and $C$ are the signed boundary formula and the
shared-child identity. The remaining equality is a direct substitution in the
three cases $j=2$, $j=3$, and $j\ge4$: the constant suffix of $D$ determines
its common child $X$, while the one-step longer factor state $F$ has the same
raw side $Y$. For the outcome assertion, a common child of a \DRAW\ cannot be
\LOSS. If the exponent-one member is \DRAW, the frame $\{D,C\}$ already
contains a \DRAW. In the parent form of an obligation, the only remaining
orientation has the exponent-two member \DRAW; if $C$ were \WIN, outcome
recursion through the displayed diamond would force $Y$ both \LOSS\ and a
child of the \DRAW\ state $F$, impossible. Hence $C$ is \DRAW. For $j=2$,
$D,C$ are exactly the exponent-two/one pair over $J(y)$.
\end{proof}


\section{Audit-hardened routing reduction}\label{sec:routing}

The preceding sections are local arithmetic and outcome lemmas.  The remaining
issue is not another binary identity; it is the global bookkeeping needed to
compose the local diamonds without silently installing an unrelated proof
height or silently changing a ranked source anchor.

\subsection{The old anchor shortcut is permanently withdrawn}

The earlier routing proof attached a retained integer $p$ to an $A$-streak and,
at a later $B$-selecting cursor $y$, used an estimate of the form
\[
 y\le \frac{3p+1}{2}
\]
to conclude that a high-valuation $B$ step lowered the retained source.  That
provenance is false after more than one canonical $A$ continuation.

\begin{lemma}[Anchor stress test]\label{lem:anchor-stress}
Let $x_0=20$ and $x_{i+1}=A(x_i)$.  Then
\[
 x_0,x_1,x_2,x_3=20,30,45,68.
\]
The canonical phases remain $A$-selecting through the first three sources and
switch at $68$.  The selected $B$ transition at $68$ has valuation $3$ and
\[
 B(68)=25>20.
\]
Consequently no proof may infer a decrease relative to the original anchor
$20$ merely from the later high-valuation $B$ transfer.
\end{lemma}

\begin{proof}
Directly,
\[
 A(20)=30,\qquad A(30)=45,\qquad A(45)=68,
\]
and
\[
 \alpha(20),\alpha(30),\alpha(45),\alpha(68)=1,0,1,1.
\]
The lifted phase flips at every canonical $A$ continuation, so the required
phase at $68$ is $0$, opposite to $\alpha(68)=1$.  Moreover
\[
 3J(68)+1=616=2^3\cdot77,
 \qquad 77=J(25),
\]
which gives the asserted valuation and $B(68)=25$.
\end{proof}

The revised routing architecture therefore treats the displayed integers in an
$A$-streak as \emph{temporary cursors}.  A cursor is never promoted to a ranked
source merely because it is smaller than a later cursor.  A control reset may
be justified only by one of the two certificates below:
\begin{enumerate}[label=(R\arabic*)]
\item an actual DRAW state whose retained coefficient source is proved smaller
than the currently ranked source by an exact factor-free/source lemma;
\item a finite WIN/LOSS proof token that is proved to be a proper proof-tree
descendant of a token already present in the ranked multiset (or the unique
initial token installation).
\end{enumerate}

\subsection{Finite control is not semantic completeness}

The finite control vocabulary consists of the following names:
\begin{center}
\begin{tabular}{lll}
\texttt{boundary\_entry} & \texttt{loss\_sibling\_entry} & \texttt{a\_obligation}\\
\texttt{a\_test\_4} & \texttt{a\_test\_3} & \texttt{a\_test\_2}\\
\texttt{a\_test\_1} & \texttt{b\_select} & \texttt{b2\_first}\\
\texttt{b2\_ready} & \texttt{factor\_fork} & \texttt{high\_return}\\
\texttt{marked\_tail} & \texttt{short\_lift} & \texttt{terminal\_macro}.
\end{tabular}
\end{center}
The finite graph of these names is only a bookkeeping device.  A graph edge is
admissible in the mathematical proof only when its guard is derived from the
actual game states and its finite witnesses have the exact parent/child
incidence claimed for them.  The accompanying JSON file therefore calls its
edges \emph{declared routing edges}, not proved game transitions.

The arithmetic part of a $B$-selecting transfer is safe without any retained
anchor comparison.  Lemma~\ref{lem:b-diamond} gives an adjacent frame over the
actual selected source $z=B(y)$.  If the selected valuation is $2$, this frame
is again an exact obligation.  If the valuation is at least $3$, it is routed
to the factor-frame normalizer.  In neither row do we assert that $z$ is below
an older pre-streak source.

\section{The two remaining proof obligations}\label{sec:remaining}

The following statements are deliberately separated from the proved lemmas.
They are the only assertions needed to turn the present reduction back into an
unconditional no-DRAW proof, and they are \emph{not} claimed proved in this
version.

\paragraph{Obligation P1: seeded $A/B_2$ reset provenance.}
Suppose the inner routing layer already contains a ranked finite proof token.
Follow an actual DRAW through any finite canonical $A$-streak and then through
a valuation-two $B$ transfer that returns to an obligation.  Before the
control is reset to a new \texttt{a\_obligation}, one must exhibit either
(i) a genuine ranked source decrease, or (ii) a finite witness which is a
strict proof-tree descendant of an already ranked token.  The statement must
include multi-$A$ streaks and all four exceptional residue classes.  It is not
enough that the new source is WIN, and it is not enough that it is smaller
than the immediate cursor.

\paragraph{Obligation P2: semantic exhaustiveness of the factor router.}
For every actual DRAW entering any factor-router control, the complete
outcome split must be derived from the game relation.  The relevant controls
are \texttt{factor\_fork}, \texttt{high\_return},
\texttt{marked\_tail}, \texttt{short\_lift}, and
\texttt{terminal\_macro}.  Every case must land either in a declared
non-reset routing edge or in a strict source/token edge.  In particular, every reset to a higher control
must carry the explicit descendant certificate for the token it replaces.
No finite endpoint of unrelated height may be inserted merely because it is
WIN or LOSS.

The long proof ledger in the repository contains many of the local identities
needed for P1 and P2, including the factor-fork and high-return descendants.
They are useful evidence, but the present manuscript does not promote the two
global composition statements to theorems until their guards and provenance
are checked as one closed transition system.

\begin{proposition}[Conditional routing closure]\label{prop:conditional-routing}
Assume P1 and P2.  Then the marked routing relation admits a well-founded
lexicographic/multiset ranking in which every control reset is preceded by a
strict retained source or certified proof-token decrease, while all remaining
routing is finite by the bounded control graph and its explicit tail/factor
counters.
\end{proposition}

\begin{proof}
The proof is the standard well-founded-fibre assembly.  Rank finite proof
witnesses by the natural multiset sum
\[
 \Phi(M)=\mathop{\#}_{h\in M}\omega^h
\]
from Lemma~\ref{lem:token-rank}.  A certified descendant replacement strictly
lowers $\Phi$.  A genuine retained-source edge lowers the numerical source and
may reset later data.  At fixed retained source and token multiset, P1 and P2
say exactly that no edge which resets the finite control is unranked.  Delete
the strict source/token edges.  What remains is the declared finite control
graph together with same-control recurrences whose explicit natural tail or
factor counter decreases.  The residual relation is well-founded.  The
well-founded-fibre lemma then gives the full marked relation.
\end{proof}

\section{Conditional no-DRAW consequence and present status}

\begin{theorem}[Conditional no-DRAW theorem]\label{thm:conditional}
If Obligations P1 and P2 hold, then the target no-DRAW statement of
Section~\ref{target:nodraw} follows.
\end{theorem}

\begin{proof}
Assume a DRAW exists and choose a least coefficient source among DRAW
positions.  The long-tail reduction, the four explicit exponent-one/two
boundary cases of Proposition~\ref{prop:boundary-reduction}, and the
factor-free base entry of Proposition~\ref{prop:base-entry} yield either a
strict smaller-source DRAW (contradicting minimality) or a typed entry.  The
factorful exponent-one row is handled by Proposition~\ref{prop:factorful-entry};
no reverse-factor lemma is used at exponent one and no coefficient/source
shortcut is used with a factor $3^k$.

Under P1 and P2, Proposition~\ref{prop:conditional-routing} makes every
outcome-compatible marked continuation well-founded.  But every DRAW has no
LOSS child and at least one DRAW child, so the local outcome diamonds could be
iterated forever while retaining an actual DRAW member.  This contradicts
well-foundedness.  Hence no DRAW exists in the conjugated game, and
Proposition~\ref{prop:first-normal} transfers finite remoteness to every
original positive odd start.
\end{proof}

\paragraph{Status.}
The implication above is proved; P1 and P2 are not.  Accordingly version 5.0
\textbf{does not claim that the prize problem is solved}.  This wording is
intentional.  It prevents a finite control certificate, a regression test, or
a moving repository supplement from being mistaken for the missing semantic
proof.

\section{Verification boundary and reproducibility}

The package contains three kinds of checks.
\begin{itemize}[leftmargin=*]
\item \texttt{verify\_v5\_0.py} checks the exact arithmetic identities used in
the proved part, the old source/coefficient counterexample, the permanent
$20\to30\to45\to68$ anchor counterexample, and the declared finite control
graph.
\item \texttt{routing\_certificate\_v5\_0.json} records controls and guards.
Every edge is labelled either \texttt{pure-routing}, \texttt{strict-source},
\texttt{strict-token}, or \texttt{OPEN-P1/P2}.  The verifier rejects an
unlabelled reset edge.
\item \texttt{OPEN\_PROOF\_OBLIGATIONS\_v5\_0.md} gives the exact mathematical
targets still required for an unconditional theorem.
\end{itemize}

These checks are deliberately narrower than a formal proof.  No end-to-end
Lean formalization is claimed.

\appendix
\section{Permanent adversarial regression cases}\label{app:regressions}

The following two examples are normative negative tests for every future
revision.

\paragraph{Coefficient/source separation.}
For $s=1$, $J(1)=5$, $a=3J(1)=15$, and $e=1$, the common boundary child is
$R(44)=22$.  Its canonical coefficient is $11<15$, but $11=J(3)$, so the
coefficient source grows from $1$ to $3$.  Therefore a smaller canonical
coefficient does not imply a smaller source when powers of three are present.

\paragraph{Pre-streak anchor separation.}
The calculation of Lemma~\ref{lem:anchor-stress} gives
$20\to30\to45\to68$ and $B(68)=25>20$.  Therefore a high-valuation $B$
transfer after several $A$ steps does not by itself lower the source retained
at the beginning of the streak.

\section{Declared routing inventory}\label{app:routing}

The table is a checklist, not a proof of P1/P2.
\begin{center}
\small
\begin{tabular}{@{}p{.25\linewidth}p{.61\linewidth}@{}}
\toprule
Control & Allowed non-strict successor classes \\
\midrule
\texttt{boundary\_entry} & finite tail countdown; then
  \texttt{a\_obligation} or \texttt{factor\_fork} \\
\texttt{loss\_sibling\_entry} & finite tail countdown; then
  \texttt{factor\_fork} \\
\texttt{a\_obligation} & bounded $A$ test or factor fork; a smaller-source
  DRAW may re-enter only through a proved source comparison \\
\texttt{a\_test\_4,3,2,1} & next lower test or \texttt{b\_select} \\
\texttt{b\_select} & valuation $2$ to \texttt{b2\_first}; valuation
  $\ge3$ to \texttt{factor\_fork}; no retained-anchor descent is asserted \\
\texttt{b2\_first} & \texttt{b2\_ready} \\
\texttt{b2\_ready} & reset requires P1 \\
\texttt{factor\_fork} & \texttt{high\_return} or strict certified exit \\
\texttt{high\_return} & reset requires the P2 descendant certificate \\
\texttt{marked\_tail} & shorter tail/short lift, or strict P2 token exit \\
\texttt{short\_lift} & \texttt{a\_obligation} \\
\texttt{terminal\_macro} & every reset requires a P2 token certificate \\
\bottomrule
\end{tabular}
\end{center}

\begin{thebibliography}{9}
\bibitem{Guy1986}
R. K. Guy, John Isbell's Game of Beanstalk and John Conway's Game of
Beans-Don't-Talk, \emph{Mathematics Magazine} 59(5) (1986), 259--269.

\bibitem{Guy1996}
R. K. Guy, Unsolved Problems in Combinatorial Games, Problem 42, in
\emph{Games of No Chance}, MSRI Publications 29, Cambridge University Press,
1996.

\bibitem{Althofer2024}
I. Alth\"ofer, M. Hartisch, T. Zipproth, Analysis of a Collatz Game and Other
Variants of the $3n+1$ Problem, \emph{Advances in Computer Games} (2024),
123--132.

\bibitem{AlthoferPage}
I. Alth\"ofer, Collatz Problems: The Two-Player $3n+-1$ Game,
\url{https://althofer.de/collatz-prizes.html}.

\bibitem{Repo}
G. Pochuev, Optimal $3n\pm1$ Game: Proof and Verification Repository,
\url{https://github.com/Grisha-Pochuev/3n-plus-minus-1-game}.
\end{thebibliography}

\end{document}
